\subsection{Set Builder Notation} 
If we have a set that's hard to describe in roster notation, we use an alternate notation.
\begin{definition}[Set-Builder Notation]
A set written in \term{set-builder notation} has the following form:
\[
\left\{\mbox{form of a set element} \mid \mbox{rules variables must follow}\right\}
\]
\end{definition}
\begin{Example}
The set
\[
\left\{
\frac{n}{2}
\mid
n\mbox{ is a natural number less than 10}
\right\}
\]
is the set
\[
\makeblank[3in]
\]
\end{Example}
\begin{Note}
We read this as ``the set of all (elements) such that (rule).'' For example, $\left\{x\mid x<10\right\}$ would be read as:\makeblanknew
\end{Note}
